\section{Dictionary and remarks on the original preprint}\label{app:dictionary}

\paragraph{Correspondence of results.}
\begin{center}\small
\begin{tabular}{@{}ll@{}}
\toprule
This document & Original preprint\\
\midrule
\cref{thm:main} & Theorem 1.1\\
\cref{lem:strong,lem:subset-deficit} & Lemma 2.2\\
\cref{lem:internal,lem:blowup} & Lemma 2.3\\
\cref{prop:reduction} & Proposition 2.1\\
\cref{lem:pruning} & Lemma 3.1\\
\cref{lem:kernel} & Lemma 3.2\\
\cref{prop:matching-bound} & the matching bound inside the proof of Lemma 3.1\\
\cref{prop:no-concave} & Proposition 4.1\\
\cref{cor:weighted,cor:triangles} & Corollaries 5.1 and 6.1\\
\bottomrule
\end{tabular}
\end{center}

\paragraph{Changes of notation.}
To avoid clashes we write $\mu=\abs P+\abs Q$ (the original uses $M$, which
we reserve for matchings $\Mm$ and the matrix of \cref{cor:weighted});
$\unc(P,Q)$ for the number $d$ of uncovered points; $\Phi^{p,i}$ for the
matrix $D^{p,i}$ (since $D$ denotes a graph); $\overline N(v)$ for the
non-neighbors $Z(v)$ (since $Z$ denotes conflict points); and
$\mathrm{Arc}(D)$ for the arc set $A(D)$ (since $A$ is a matrix). We added the
names $\val$, ``admissible'', $\Gamma_L$, $\Gamma_N$, $R_1$, $\kappa$.
Otherwise the notation is that of the original.

\paragraph{Remarks.}
We checked every step of the original and found \textbf{no mathematical
error}. The following points are only implicit there, and we have made them
explicit.
\begin{enumerate}[leftmargin=*]
\item Lemma~2.2 in the original states strong connectivity together with
the subset deficit. Strong connectivity is exactly what is needed for
$E=F(S)\setminus S\ne\varnothing$, and it is the reason why the number of
vertices is minimized before the number of arcs (\cref{lem:strong}).
\item The proof of the blow-up lemma uses only the integrality of the
subset deficit and the inequality~\eqref{eq:blowup-master}. The latter holds
in \emph{every} blow-up of every oriented graph. We isolated it, together
with its unblown-up version~\eqref{eq:error-term}, to explain why the
blow-up is used (\cref{rem:error-term}).
\item In the pruning lemma, the bound~\eqref{eq:pruning-bound} is not a
bound on the minimum number of deletions (\cref{ex:K22}). This is easy to
misread, so we added an explicit example.
\item In the matching bound, the step $\rank(NL)\ge\ell+t-\abs Z$ is Sylvester's
rank inequality, and the choice of a maximum matching with the fewest
conflict edges is used only to make $\alpha,\beta$ maximum in their
subgraphs (\cref{cl:alpha-max}).
\item Seacrest's equivalences~\cite{Seacrest2015} are needed only for the
arc-weighted form. The vertex-weighted forms follow from \cref{thm:main}
by the short argument in \cref{sec:weighted}.
\item The bibliographic and historical statements in the introduction are
taken from the original preprint and were not independently verified.
\end{enumerate}

\section{Computer checks and formalization}\label{app:checks}

\sloppy
\paragraph{Randomized tests.}
The folder \texttt{checks/} contains three dependency-free Python scripts.
\texttt{check\_pruning.py} builds, for thousands of random relations on
$2$--$5$ points and random families $\Rf,\Cf$, the graph $\Gamma$. It verifies
$\nu(\Gamma)\le\abs Z+\abs H$ (\cref{prop:matching-bound}), checks that the
K\H{o}nig cover produces a pruning with \ref{P1}--\ref{P3}, and confirms by
exhaustive search that a valid pruning exists. The bound is attained in
roughly a quarter of the instances with conflicts, so it is sharp.
\texttt{check\_reduction.py} verifies the unconditional
inequality~\eqref{eq:blowup-master} for random blow-ups, the increment
identity~\eqref{eq:g-increment}, and \cref{cl:diag-free,cl:H-uncovered} for
random admissible pairs (these do not need concavity).
\texttt{check\_small\_graphs.py} runs through \emph{all} labelled oriented
graphs on at most $6$ vertices ($14{,}348{,}907$ on six vertices; about two
minutes). It confirms \cref{thm:main} for each of them and confirms
\cref{prop:no-concave} directly: none of the $5{,}545{,}708$ graphs on six
vertices with positive indegrees is strictly concave. All tests pass. They are sanity checks, not proofs.

\paragraph{Lean formalization.}
The folder \texttt{SeymourLean/} is a Lean~4 project using Mathlib (tag
\texttt{v4.35.0}) that formalizes the \emph{entire} proof of
\cref{thm:main}. The main formal statement is
\begin{center}\small
\texttt{theorem Seymour.seymour\_second\_neighbourhood (V : Type) [Fintype V]
[DecidableEq V] [Nonempty V] (adj : V $\to$ V $\to$ Prop) [DecidableRel adj]
(hor : IsOriented adj) : $\exists$ v, (N1 adj v).card $\le$ (N2 adj v).card}
\end{center}
where \texttt{IsOriented} means irreflexive and asymmetric, \texttt{N1 v} is
the set of out-neighbors, and \texttt{N2 v} is the set of $w\ne v$ with
$v\not\to w$ and $v\to u\to w$ for some $u$. \texttt{lake build} checks it with
no \texttt{sorry}, and \texttt{\#print axioms} reports only Lean's standard
axioms \texttt{propext}, \texttt{Classical.choice} and \texttt{Quot.sound}.
The file \texttt{SeymourLean/build.log} records a full build together with
the output of \texttt{lake env lean CheckAxioms.lean}, so this can be checked
without rebuilding. The files follow the paper:
\begin{itemize}
\item \texttt{Reduction.lean}: minimal counterexample, \cref{lem:strong} and
\cref{lem:subset-deficit};
\item \texttt{BlowUp.lean}: \cref{lem:blowup};
\item \texttt{Contradiction.lean}: \cref{prop:no-concave} from the pruning lemma;
\item \texttt{Konig.lean}, \texttt{Gamma.lean}: K\H{o}nig's theorem (derived from
Mathlib's Hall theorem) and \cref{cl:cover-to-pruning};
\item \texttt{LinAlg.lean}, \texttt{Generic.lean}, \texttt{MatchingCore.lean},
\texttt{MatchingBound.lean}: Sylvester's inequality, the kernel lemma
(\cref{lem:kernel}), $NL=AB$, \cref{cl:conflict-kernel},
\cref{prop:matching-bound}, and \cref{lem:pruning};
\item \texttt{Final.lean}: the unconditional theorem;
\item \texttt{Insights.lean}: the remarks of the insights pass, see below.
\end{itemize}
The remarks added in \cref{sec:story,sec:ledger,sec:pruning,sec:contradiction}
are formalized in \texttt{Insights.lean}: \texttt{F1\_F2\_comm} ($\Ff\Fs=\Fs\Ff$)
and \texttt{kronecker\_comm} (\cref{rem:kronecker});
\texttt{diag\_compatible\_iff\_oriented} and
\texttt{exists\_admissible\_iff\_oriented} (\cref{lem:admissible-iff-oriented});
\texttt{conflict\_swap}, \texttt{Zset\_swap}, \texttt{Hset\_swap}
(\cref{rem:reversal}); \texttt{gain\_identity} (the exact form of
\cref{cl:gain}) and \texttt{moved\_identity} (\cref{rem:move-repair});
and \texttt{step\_ledger}, \texttt{step\_value}: for an arbitrary relation
and any admissible pair, pruning the complements gives an admissible pair of
value at least $\val(P,Q)-\sum\kappa$ (\cref{rem:algorithm}), with no
concavity assumed. Not formalized: the Dulmage--Mendelsohn reading of the
kernel lemma (\cref{rem:DM}), the computer observations of
\cref{rem:combinatorial}, and the running-time statement of
\cref{rem:algorithm}. \texttt{CheckAxioms.lean} prints the axioms of
these theorems as well.
A few steps are argued differently in Lean. Generic coefficients are taken
as the variables themselves in the fraction field of a polynomial ring
rather than as a real point. Only $\rank L\ge\ell$ is used, obtained from
Hall's theorem and an exchange argument, instead of the full statement
$\rank L=\nu(\Gamma_L)$. The rank count goes through the three-term complex of \cref{sec:rank}
(\texttt{ComplexCount.lean}: \texttt{complex\_count} is the counting of
\cref{cl:kerA,cl:rank-inequality}, and \texttt{dimension\_count} is
\cref{cl:rank-inequality} for the matrices of the paper); the earlier
version of the project used Sylvester's inequality twice instead. Strong connectivity is replaced by ``every nonempty
proper set has an arc leaving it''. The formalization found no mathematical
error in the paper. It also shows that the pruning lemma holds for an
arbitrary relation and that orientation is used exactly at the places listed
in \cref{rem:hypotheses}.

